\usepackage{microtype}
\usepackage{graphicx}
%\usepackage{subfigure}
\usepackage{booktabs} % for professional tables
\usepackage{hyperref}


\usepackage[utf8]{inputenc} % allow utf-8 input
\usepackage{url}            % simple URL typesetting
%\usepackage{amsfonts}       % blackboard math symbols
\usepackage{nicefrac}       % compact symbols for 1/2, etc.

% additional packages
\usepackage[british]{babel}
\usepackage{natbib}
\usepackage{enumerate}
\usepackage{cancel}
% \usepackage{xcolor}
\usepackage{soul}
\usepackage{multirow}
\usepackage{mathtools}
\usepackage{amsmath,amsthm,verbatim,amssymb,amsfonts,amscd}
\usepackage{commath}
\usepackage{bbm}
\usepackage{thmtools,thm-restate}
\usepackage{cleveref} % load last

\babelhyphenation[british]{ap-pro-xi-ma-te}
\babelhyphenation[british]{em-pi-ri-cal}
\babelhyphenation[british]{re-lia-bi-li-ty}
\babelhyphenation[british]{va-li-da-tion}
\babelhyphenation[british]{me-chan-ism}
\babelhyphenation[british]{si-mul-ta-ne-ous-ly}


% allows equations to be split over mutliple pages
\allowdisplaybreaks

% math statements
\newtheorem{theorem}{Theorem}
\newtheorem{lemma}[theorem]{Lemma}
\newtheorem{proposition}[theorem]{Proposition}
\newtheorem{remark}[theorem]{Remark}
\newtheorem{corollary}[theorem]{Corollary}
\newtheorem{definition}[theorem]{Definition}

% custom commands
\DeclareMathOperator*{\argmin}{\arg\!\min}
\DeclareMathOperator*{\argmax}{\arg\!\max}
\DeclareMathOperator*{\E}{\mathbb{E}}
\DeclareMathOperator*{\Var}{\mathbb{V}}
\DeclareMathOperator*{\Covar}{Cov}
\DeclareMathOperator{\Prob}{\mathbb{P}}
\DeclareMathOperator{\diverg}{D}
\DeclareMathOperator{\law}{Law}
\DeclareMathOperator{\lipschitz}{Lip}
\DeclareMathOperator{\poly}{poly}
\DeclareMathOperator{\diam}{diam}
\DeclareMathOperator{\sign}{sign}
\DeclareMathOperator{\proj}{proj}
\DeclareMathOperator{\vectorise}{vect}
\DeclareMathOperator{\supp}{supp}
\DeclareMathOperator{\given}{\vert}
\DeclareMathOperator{\conv}{\star}
\DeclareMathOperator{\diag}{diag}
% \DeclareMathOperator{\KL}{\mathrm{KL} \, }
\DeclarePairedDelimiter\floor{\lfloor}{\rfloor}
\DeclarePairedDelimiter{\parenth}{\lparen}{\rparen}
\DeclarePairedDelimiter{\bracket}{\lbrack}{\rbrack}
\DeclarePairedDelimiter{\vertbars}{\vert}{\vert}
\DeclarePairedDelimiterX{\infdivx}[2]{(}{)}{%
	#1 \, \delimsize\| \, #2%
}
\DeclarePairedDelimiterX{\commaparen}[2]{(}{)}{%
	#1 , #2%
}
\DeclarePairedDelimiterX{\innerprod}[2]{\langle}{\rangle}{%
	#1 , #2%
}
\DeclarePairedDelimiterX{\parenththree}[3]{\lparen}{\rparen}{
	#1 ; #2 ; #3	
}
\newcommand{\niton}{\not\mathrel{\text{\reflectbox{$\in$}}}}
\newcommand*\xbar[1]{%
	\hbox{%
		\vbox{%
			\hrule height 0.5pt % The actual bar
			\kern0.5ex%         % Distance between bar and symbol
			\hbox{%
				\kern-0.1em%      % Shortening on the left side
				\ensuremath{#1}%
				\kern-0.1em%      % Shortening on the right side
			}%
		}%
	}%
} 

\newcommand{\fun}[1]{#1}
\newcommand{\enorm}[1]{\norm{#1}_2}
%\newcommand{\abs}{\vertbars}
\newcommand{\KL}{\mathrm{KL} \, \infdivx}
\newcommand{\TV}{\mathrm{TV} \commaparen}
\newcommand{\expect}[1]{\E_{#1} \parenth}
\newcommand{\divg}{\diverg \, \infdivx}
%\newcommand{\inftsm}{\mathrm{d} \hspace{1pt} \parenth}
\newcommand{\inftsm}[1]{\mathrm{d} #1}
\newcommand{\lebesgue}{\lambda}
\newcommand{\gradwrt}[1]{\nabla_{#1}}
%\newcommand{\proj}{\mathrm{proj} \, \parenth}
\newcommand{\sigmaAlg}{\sigma \, \parenth}
%\newcommand{\borelAlg}{\mathcal{B} \, \parenth}
\newcommand{\entropy}{\mathcal{H} \, \parenth}
\newcommand{\order}{\mathcal{O} \, \parenth}
\newcommand{\determ}{\vertbars}
\newcommand{\vect}[1]{\boldsymbol{#1}}
\newcommand{\mat}[1]{\boldsymbol{#1}}
\newcommand{\aset}[1]{#1}
\newcommand{\aspace}[1]{\mathcal{#1}}
\newcommand{\rv}[1]{\expandafter\MakeUppercase\expandafter{#1}}
\newcommand{\distr}[1]{\expandafter\MakeUppercase\expandafter{\mathrm{#1}}}
%\newcommand{\distr}[1]{\mathrm{#1}}
\newcommand{\dens}[1]{#1}
\newcommand{\const}[1]{\mathrm{#1}}
\newcommand{\complem}[1]{#1^\mathrm{C}}
\newcommand{\transp}[1]{#1^\mathrm{T}}
\newcommand{\inv}[1]{#1^{-1}}
\newcommand{\pseudoinv}[1]{#1^{\dag}}
\newcommand{\trace}[1]{\mathrm{Tr}\left( #1 \right)}
\newcommand{\natnum}{\mathbb{N}}
\newcommand{\hilbert}[1]{\mathcal{#1}}
\newcommand{\dirac}{\fun{\delta}}
\newcommand{\diracf}{\dirac \, \parenth}
\newcommand{\kronecker}{\fun{\delta}_{\mathrm{Kr}}}
\newcommand{\kroneckerf}{\kronecker \, \parenth}
% \newcommand{\indicator}[1]{\mathbb{I}_{#1}}
\newcommand{\indicator}[1]{\mathbbm{1}_{#1}}
\newcommand{\indicatorf}[2]{\indicator{#1} ( #2 )}
\newcommand{\identity}{\mathrm{Id}}
\newcommand{\identityf}[1]{\identity(#1)}
\newcommand{\e}[1]{\const{e}^{#1}}
\newcommand{\factorial}[1]{#1!}


% \newcommand{\kernel}{\mathcal{K}}
\newcommand{\kernel}{\kappa}
\newcommand{\kernhat}{\hat{\kernel}}
\newcommand{\kerntilde}{\tilde{\kernel}}
\newcommand{\kernelf}[2]{\kernel_{#1}^{#2} \, \commaparen}
\newcommand{\kernhatf}[2]{\kernhat_{#1}^{#2} \, \commaparen}
\newcommand{\kerntildef}[2]{\kerntilde_{#1}^{#2} \, \commaparen}
\newcommand{\ntk}{\Theta}
\newcommand{\ntktilde}{\widetilde{\ntk}}
\newcommand{\ntkhat}{\widehat{\ntk}}
\newcommand{\ntkf}[2]{\ntk_{#1}^{#2} \, \commaparen}
\newcommand{\ntktildef}[2]{\ntktilde_{#1}^{#2} \, \commaparen}
\newcommand{\ntkhatf}[2]{\ntkhat_{#1}^{#2} \, \commaparen}
\newcommand{\interpKern}{\mathcal{I}}
% \newcommand{\interpRand}{\interpKern^{r}}
% \newcommand{\interpStruct}{\interpKern^{s}}
\newcommand{\interpKernCoeff}{\alpha}

\newcommand{\R}[1]{\mathbb{R}^{#1}}
\newcommand{\Q}[1]{\mathbb{Q}^{#1}}
\newcommand{\Lp}[1]{\mathrm{L}^{#1}}
\newcommand{\contFns}[1]{C(#1)}
\newcommand{\contBoundFns}[1]{C_b(#1)}
\newcommand{\contVanishFns}[1]{C_0(#1)}

\newcommand*\rfrac[2]{{}^{#1}\!/_{#2}}

\newcommand{\gauss}{\mathcal{N}}
\newcommand{\gaussDens}{\fun{\phi}}
\newcommand{\gaussDensFn}[2]{\gaussDens_{#1, #2}}
\newcommand{\standGaussDensFn}{\gaussDensFn{0}{1}}
\newcommand{\mgauss}{\mathcal{MN}}
\newcommand{\gp}{\mathcal{GP}}
\newcommand{\categorical}{\text{Categorical}}
\newcommand{\bernoulli}{\text{Bernoulli}}
\newcommand{\chisq}{\chi^2}

% \newcommand{\convergeDist}{\overset{\mathcal{L}}{\longrightarrow}}
\newcommand{\convergeDist}{\rightsquigarrow}
\newcommand{\convergeProb}{\overset{P}{\longrightarrow}}
\newcommand{\convergeAs}{\overset{a.s.}{\longrightarrow}}

% function restriction
\newcommand\restr[2]{{% we make the whole thing an ordinary symbol
  \left.\kern-\nulldelimiterspace % automatically resize the bar with \right
  #1 % the function
%   \vphantom{\big|} % pretend it's a little taller at normal size
  \right|_{#2} % this is the delimiter
  }}


  

\newcommand{\generalMean}{\bar{\generalSum}}

\newcommand{\genericDimenstion}{d}
\newcommand{\inputDimension}{\genericDimenstion^{0}}
\newcommand{\outputDimension}{\genericDimenstion^{\depth + 1}}
\newcommand{\spatialDimension}{\genericDimenstion^s}
\newcommand{\layerDimension}[1]{\genericDimenstion_{\sequenceVariable}^{#1}}
% \newcommand{\layerDimensionN}{\genericDimenstion_{\sequenceVariable}^{\depthSymbol}}
\newcommand{\layerDimensionN}{\layerDimension{\depthSymbol}}


\newcommand{\headIndex}{h}
\newcommand{\headSymbol}{H}
\newcommand{\headN}{\headSymbol_{\sequenceVariable}^{\depthSymbol}}

\newcommand{\softmax}{\zeta}
\newcommand{\softmaxMean}{\bar{\softmax}}
\newcommand{\params}{\theta}

\newcommand{\keySymbol}{K}
\newcommand{\querySymbol}{Q}
\newcommand{\valueSymbol}{V}
\newcommand{\logitSymbol}{G}
\newcommand{\outputSymbol}{O}
\newcommand{\key}[2]{\keySymbol_{#1}^{#2}}
\newcommand{\query}[2]{\querySymbol_{#1}^{#2}}
\newcommand{\val}[2]{\valueSymbol_{#1}^{#2}}
\newcommand{\logit}[2]{\logitSymbol_{#1}^{#2}}
\newcommand{\keyN}{\key{\sequenceVariable}{\depthSymbol \headIndex}}
\newcommand{\queryN}{\query{\sequenceVariable}{\depthSymbol \headIndex}}
\newcommand{\valueN}{\val{\sequenceVariable}{\depthSymbol \headIndex}}
\newcommand{\logitN}{\logit{\sequenceVariable}{\depthSymbol \headIndex}}

\newcommand{\posEmbSymbol}{E}
\newcommand{\posEmb}[2]{\posEmbSymbol_{#1}^{#2}}
\newcommand{\posEmbN}{\posEmb{\sequenceVariable}{\depthSymbol \headIndex}}
\newcommand{\covPosEmb}[2]{R_{#1}^{#2}}


\newcommand{\weightMatSymbol}{W}
\newcommand{\W}{\weightMatSymbol}
\newcommand{\WQ}{\weightMatSymbol^{\querySymbol}}
\newcommand{\WK}{\weightMatSymbol^{\keySymbol}}
\newcommand{\WV}{\weightMatSymbol^{\valueSymbol}}
\newcommand{\WO}{\weightMatSymbol^{\outputSymbol}}
\newcommand{\weightQ}[1]{\weightMatSymbol_{\sequenceVariable, #1}^{\depthSymbol \headIndex, \querySymbol}}
\newcommand{\weightK}[1]{\weightMatSymbol_{\sequenceVariable, #1}^{\depthSymbol \headIndex, \keySymbol}}
\newcommand{\weightV}[1]{\weightMatSymbol_{\sequenceVariable, #1}^{\depthSymbol \headIndex, \valueSymbol}}
% \newcommand{\weightO}[1]{\weightMatSymbol_{\sequenceVariable, #1}^{\depthSymbol, \outputSymbol}}
\newcommand{\weightOGen}[2]{\weightMatSymbol_{#1}^{#2}}
\newcommand{\weightO}[1]{\weightOGen{\sequenceVariable, #1}{\depthSymbol, \outputSymbol}}
\newcommand{\weightsQ}{\weightMatSymbol_{\sequenceVariable}^{\depthSymbol \headIndex, \querySymbol}}
\newcommand{\weightsK}{\weightMatSymbol_{\sequenceVariable}^{\depthSymbol \headIndex, \keySymbol}}
\newcommand{\weightsV}{\weightMatSymbol_{\sequenceVariable}^{\depthSymbol \headIndex, \valueSymbol}}
\newcommand{\weightsO}{\weightMatSymbol_{\sequenceVariable}^{\depthSymbol, \outputSymbol}}

% \newcommand{\headDimension}{\genericDimenstion_{\sequenceVariable}^{\headSymbol^{\depthSymbol}}}
\newcommand{\headDimension}{\genericDimenstion_{\sequenceVariable}^{\depthSymbol, \headSymbol}}
% \newcommand{\headDimension}{\headSymbol_{\sequenceVariable}^{\depthSymbol}}
\newcommand{\keyDimension}{\genericDimenstion_{\sequenceVariable}^{\depthSymbol, \keySymbol}}
\newcommand{\queryDimension}{\genericDimenstion_{\sequenceVariable}^{\depthSymbol, \querySymbol}}
\newcommand{\valueDimension}{\genericDimenstion_{\sequenceVariable}^{\depthSymbol, \valueSymbol}}
\newcommand{\logitDimension}{\genericDimenstion_{\sequenceVariable}^{\depthSymbol, \logitSymbol}}
\newcommand{\peDimension}{\genericDimenstion_{\sequenceVariable}^{\depthSymbol, \posEmb{}{}}}


\newcommand{\stdSymbol}{\sigma}
\newcommand{\keyStd}{\stdSymbol_{\keySymbol}}
\newcommand{\queryStd}{\stdSymbol_{\querySymbol}}
\newcommand{\valueStd}{\stdSymbol_{\valueSymbol}}
\newcommand{\outStd}{\stdSymbol_{\outputSymbol}}
\newcommand{\QKStd}{\stdSymbol_{\querySymbol \keySymbol}}
\newcommand{\OVStd}{\stdSymbol_{\outputSymbol \valueSymbol}}
\newcommand{\keyVar}{\keyStd^2}
\newcommand{\queryVar}{\queryStd^2}
\newcommand{\valueVar}{\valueStd^2}
\newcommand{\outVar}{\outStd^2}
\newcommand{\QKVar}{\QKStd^2}
\newcommand{\OVVar}{\OVStd^2}

\newcommand{\activations}[2]{\activationSymbol_{#1}^{#2}}
\newcommand{\activationsN}{\activations{\sequenceVariable}{\depthSymbol}}
\newcommand{\channels}[2]{\tilde{\activationSymbol}_{#1}^{#2}}
\newcommand{\channelsN}{\channels{\sequenceVariable}{\depthSymbol}}
\newcommand{\logits}[2]{\logitSymbol_{#1}^{#2}}
\newcommand{\logitsN}{\logits{\sequenceVariable}{\depthSymbol}}

\newcommand{\scaling}{\tau}
\newcommand{\fIndexA}{a}
\newcommand{\fIndexB}{b}
\newcommand{\gIndexA}{a'}
\newcommand{\gIndexB}{b'}
\newcommand{\gIndexZA}{i'}
\newcommand{\gIndexZB}{j'}

\newcommand{\projectionN}{\projection_{\sequenceVariable}}
\newcommand{\summandN}{\summand_{\sequenceVariable, \headIndex}}

\newcommand{\projectionLogitSymbol}{\projection}
\newcommand{\summandLogitSymbol}{\varphi}
\newcommand{\projectionLogit}{\projectionLogitSymbol_{\sequenceVariable}^{\logitSymbol}}
\newcommand{\projectionCoefficientsLogit}{\beta}
\newcommand{\summandLogit}[1]{\summandLogitSymbol_{#1}}
\newcommand{\summandLogitN}{\summandLogit{\sequenceVariable, j}}

\newcommand{\tildeActivation}[3]{\indexedObject{\tilde{\activationSymbol}}{#1}{#2}{(#3)}}
\newcommand{\tildeActivity}[3]{\indexedObject{\tilde{\activitySymbol}}{#1}{#2}{(#3)}}
\newcommand{\tildeLogit}[2]{\widetilde{\logitSymbol}_{#1}^{#2}}
\newcommand{\tildeLogitN}[3]{\tildeLogit{#1}{#2} (#3)}
\newcommand{\uW}{\widetilde{\weightMatSymbol}}
\newcommand{\uWQ}{\uW^{\querySymbol}}
\newcommand{\uWK}{\uW^{\keySymbol}}
\newcommand{\uWV}{\uW^{\valueSymbol}}
\newcommand{\uWO}{\uW^{\outputSymbol}}
